\section{Discussion}
\label{sec:discussion}

\subsection{Error Analysis}

A rigorous error analysis is essential for understanding the failure modes of antisemitism detection systems. We categorize test-set errors into five types, based on manual review of the 50 most confident false positives and 50 most confident false negatives from our best model.

\paragraph{False positives: Political criticism flagged as antisemitic.}
The most common false positive category involves tweets that sharply criticize Israeli government policy or military actions using strong language (e.g., ``apartheid,'' ``ethnic cleansing,'' ``war crimes'') without targeting Jewish people as a group. These tweets often contain intense negative sentiment directed at a state actor, which the model conflates with antisemitic hostility. This error category is particularly consequential because incorrectly flagging legitimate political speech undermines trust in content moderation systems and can silence protected discourse.

The difficulty here is genuine: the IHRA definition acknowledges that ``criticism of Israel similar to that leveled against any other country cannot be regarded as antisemitic,'' but in practice, the line between harsh policy criticism and antisemitism-via-anti-Zionism is contested even among human annotators. Models trained on IHRA-annotated data will inevitably inherit the annotators' boundary judgments.

\paragraph{False positives: Counterspeech misclassified.}
A second false positive category involves \emph{counterspeech}: tweets that quote or reference antisemitic language in order to denounce it. For example, a tweet saying \emph{``People who call Jews `kikes' are disgusting''} is not antisemitic --- it is condemning antisemitism --- but the presence of the slur and the word ``Jews'' can trigger a positive prediction. This error type is especially relevant for the ``Kikes'' keyword group, where 65.7\% of tweets are \emph{not} antisemitic despite containing an antisemitic slur. CCI v2's heuristic symmetry filter (\S\ref{sec:cci}) avoids enforcing pair-invariance on counterspeech: tweets containing slurs are excluded from the swap manifest because the slur is identity-specific.

\paragraph{False negatives: Coded antisemitism.}
The most challenging false negatives involve \emph{coded} or \emph{implicit} antisemitism --- tweets that convey antisemitic meaning through allusion, conspiracy innuendo, or tropes without using explicit keywords. Examples include references to ``those who control the banks'' (coded reference to the antisemitic conspiracy of Jewish financial control), ``dual loyalty'' accusations, or Holocaust inversion (claiming that Jews are ``the real Nazis''). These tweets require world knowledge and cultural literacy that fine-tuned classifiers may not fully capture from limited training data. Prompted LLMs may have an advantage on coded antisemitism due to their broader world knowledge, particularly with Guided CoT prompts that explicitly list tropes to check.

\paragraph{False negatives: Sarcasm and irony.}
Ironic antisemitism --- where the speaker uses sarcasm to convey genuinely antisemitic meaning (e.g., \emph{``Of course the media is controlled by Jews, what a surprise''} said sarcastically but still propagating the trope) --- is difficult for all models. The literal content may appear as meta-commentary, but the pragmatic intent is antisemitic. Detecting such cases requires understanding ironic register, which is a known weakness of text classifiers.

\paragraph{Genuinely ambiguous cases.}
Some tweets defy clear classification because they occupy a genuine gray zone under the IHRA definition. These include tweets that criticize Israel using language that \emph{might} invoke antisemitic tropes but could also be interpreted as standard political hyperbole, or tweets that discuss Jewish identity in ways that \emph{might} essentialize but could also be read as sociological observation. Inter-annotator disagreement on such cases is expected, and model errors on them are less indicative of systematic failure.

\subsection{Why training-time invariance beats post-hoc rescaling}
\label{sec:disc_training_time}

The structural claim of \S\ref{sec:cci}, empirically confirmed in 6/6 seeds, has a precise statement: post-hoc Temperature Scaling cannot move a single counterfactual pair across the decision threshold $1/2$. The mechanism is straightforward --- monotone rescaling of logits preserves their sign --- but the implication for fairness practice is non-trivial. Temperature Scaling is the canonical calibration fix in production ML systems~\citep{guo2017calibration}, and on this task it does materially improve global ECE (from $0.077$ to $0.045$, a $-42\%$ change on Model~C, see Table~\ref{tab:main_results}). \emph{But it leaves the hard counterfactual flip rate untouched.} Practitioners who treat TS as a universal fairness fix --- and the literature implicitly does, since it is the default calibration tool in most pipelines --- are improving one fairness metric while leaving another unchanged.

The same observation applies to multicalibration (\S\ref{sec:multicalibration_tension}): although designed for subgroup-conditional calibration rather than for counterfactual fairness, it is treated in practice as a general fairness-improving post-processor. We show that on this task its effect on the L\textsubscript{1}-norm soft counterfactual flip rate is negative ($+96\%$, i.e., the rate increases). This is not a critique of multicalibration as a calibration tool; it is a critique of treating any single post-hoc tool as a universal fairness fix. CCI v2, applied at training time and explicitly designed around the pair-invariance objective, avoids the trap.

\paragraph{Why does transfer across architecture and family work?}
The structural claim above does not by itself predict that CCI v2 should transfer from a 184M-parameter DeBERTa encoder to a 440M same-family encoder and a 355M different-family encoder (RoBERTa). One might expect that the swap engine, calibrated against DeBERTa's tokenization and embedding geometry, would lose effectiveness on a different tokenizer. Empirically the opposite happens --- the cross-family CFR reduction is the largest of the three ($-73\%$ on RoBERTa-large vs.\ $-49\%$ on DeBERTa-base). A plausible explanation is that the swap engine operates at the \emph{text} level, not at the embedding level: the symmetric vocabulary is text, the tokenizer of each backbone tokenises it natively, and the loss enforces invariance over the resulting embeddings whatever they are. The intervention is therefore tokenizer-agnostic in a way that, e.g., direct embedding-space projection (LEACE-style approaches) is not.

\subsection{FNED under the corrected Dixon aggregation}
\label{sec:disc_fned}

The named pivot of this paper is CCI~v2 (fixed T, JS), which has Dixon FNED $0.793$ at base scale. Versus the strongest CLP baseline (Davani $\lambda{=}0.5$, FNED $0.786$) the gap is $\Delta = +0.007$ ($+0.9\%$, paired-bootstrap $p = 0.86$, not Bonferroni-significant). Versus same-architecture Model~C (FNED $0.797$) the gap is $\Delta = -0.004$ ($-0.5\%$, $p = 1.0$). Under the corrected aggregation the pivot is statistically indistinguishable from both comparators on FNED.

The originally pre-registered cell CCI~v2 (learned T, JS) sits higher at FNED $0.839$, a $+6.7\%$ gap to Davani $\lambda=0.5$ and $+5.3\%$ to Model~C. The corrected suite tests pivot vs.\ learned-$T$ directly: $\Delta = -0.047$, $p_{\text{boot}} = 0.27$, also not significant, so learned-$T$ is not statistically distinguishable from fixed-$T$ on FNED either, but the point estimate trend together with the within-cell ablation showing that learned-$T$ does not improve CFR\textsubscript{hard} relative to fixed-$T$ ($p = 0.25$) was the basis for our pivot change.

\paragraph{What changed under the corrected metric.}
Prior reports of the FNED Pareto trade-off and of a capacity-bound disappearance of that trade-off at large scale were computed under a range-based (max~$-$~min) aggregation that we now identify as inconsistent with the Dixon AIES 2018 definition we cite. The Dixon aggregation increases FNED magnitudes substantially (roughly $+40\%$ in absolute terms because all four groups contribute, not just the extremes) and equalises method comparisons more than the range did. The qualitative picture under the corrected metric: \emph{CCI~v2 trades nothing on FNED}, neither at base nor at large scale, against same-recipe Model~C; the FNED differences against Davani CLP variants are all within one paired-bootstrap CI of zero. The learned-$T$ cell trends worse than the fixed-$T$ pivot but the difference is no longer Bonferroni-significant.

\paragraph{Mechanism.}
CCI~v2's pair-invariance penalty pushes the model toward producing the \emph{same} probability for $(x, x')$ pairs. The learned-$T$ variant, which has an extra degree of freedom (the per-group temperature), evidently uses that freedom in a way that mildly compresses the prediction distribution at the decision boundary, slightly raising false-negative rate variance across groups; the fixed-$T$ variant has no such degree of freedom and the point-estimate FNED comes out essentially equal to Model~C's.

\paragraph{What we are not claiming.}
We do not claim CCI~v2 dominates every baseline on every fairness metric. Davani $\lambda=2.0$ has the numerically lowest Dixon FNED at base scale ($0.728$); CCI~v2 fixed-T sits in the broad mid-range of the methods we evaluate, indistinguishable from Davani $\lambda=0.5$ and Model~C. The headline claim is that CCI~v2 obtains a 49--73\% CFR reduction without any statistically significant FNED cost, which is the practitioner-relevant trade-off given the per-group n heterogeneity we describe in \S\ref{sec:fairness}.

\subsection{Multicalibration vs.\ counterfactual invariance: a tension worth surfacing}
\label{sec:disc_multicalibration}

The negative result of \S\ref{sec:multicalibration_tension} --- multicalibration post-hoc increases the L\textsubscript{1}-norm soft counterfactual flip rate by $+96\%$ on this task --- merits unpacking. Multicalibration~\citep{hebertjohnson2018multicalibration} works by partitioning the input space along protected-attribute axes and adjusting the predicted probabilities within each cell so that they match the empirical positive-class frequency. The adjustment is, in general, \emph{not} monotone: an example whose original prediction was $0.40$ may be moved to $0.55$, while another example whose original prediction was $0.55$ may be moved to $0.40$. When applied to a counterfactual pair $(x, x')$ that falls into different cells of the partition (because their identity tokens differ, by construction), the two examples can be shifted in opposite directions, increasing their probability gap even when their binary predictions agree.

This is, structurally, why CCI~v2's training-time pair-invariance does not suffer the same problem: the loss explicitly penalises \emph{within-pair} probability disagreement at the training-data level, before any cell-conditional adjustment is applied. The two notions of fairness --- subgroup-conditional calibration and counterfactual invariance --- are genuinely separate objectives, and the post-hoc multicalibration tool optimises the former at the expense of the latter on a counterfactual-rich corpus like GoldStandard2024.

The practical implication for content-moderation pipelines is that multicalibration should not be treated as a universal post-hoc fairness fix when the downstream metric of interest is counterfactual invariance. If both subgroup calibration \emph{and} counterfactual invariance are needed, training-time interventions like CCI~v2 are the appropriate path.

\subsection{Fine-tuned transformers vs.\ prompted LLMs}

Our prompted-LLM track (Mistral-7B, Qwen-2.5-7B, Llama-3.1-8B locally with 4-bit quantization, plus Llama-3.3-70B via the Groq API) is wired into the evaluation pipeline; the full multi-seed sweep is the last block of evidence pending at submission time. We outline here the expected trade-offs and revisit them in the camera-ready.

\paragraph{Accuracy and consistency.}
Fine-tuned models benefit from task-specific parameter optimization: every weight is adjusted to minimize classification error. Prompted LLMs leverage general-purpose representations and rely on the prompt to provide task specification. This typically results in higher variance for LLMs, with performance sensitive to exact prompt wording, example selection, and output parsing.

\paragraph{Reasoning transparency.}
The Guided CoT strategy (Model D3, after~\citet{patel2025evaluating}) produces step-by-step reasoning traces (target identification, trope checking, legitimate-speech assessment) that can be reviewed by human moderators. This transparency is valuable in the politically sensitive antisemitism domain. However, the reasoning may be \emph{post-hoc rather than faithful} --- the model may generate plausible reasoning that does not accurately reflect its internal decision process~\citep{turpin2023language}. Importantly, prompted LLMs do not produce calibrated probabilities in our setup; the counterfactual-stability metric we use for fine-tuned models (CFR\textsubscript{soft}) is therefore not directly applicable to D1--D3. CFR\textsubscript{hard} can be computed but requires running each LLM on both $x$ and $x'$ for the 2{,}117 swap pairs --- a cost that limits how aggressively this evaluation can scale.

\paragraph{Deployment considerations.}
Fine-tuned transformers (125--440M parameters) process a tweet in milliseconds on a GPU; the 7--8B LLMs require seconds per tweet even with 4-bit quantization, and Llama-3.3-70B over the Groq API is latency-bound by network round-trip. For platform-scale content moderation (millions of tweets per day), fine-tuned models remain the only practical option. LLMs are best suited as a second-stage classifier for ambiguous cases flagged by a fast first stage, or as an explanation generator for human moderators.

\subsection{Bimodal calibration under focal + keyword-mask training}
\label{sec:discussion:bimodality}

Cross-referencing per-epoch dev metrics with final test metrics (\S\ref{sec:bimodality}) reveals that the Model~C ECE bimodality is deterministic given the training protocol, not an inherent property of focal-loss training. Dev ECE falls monotonically with epochs for every seed --- consistent with the calibration--accuracy trade-off characterised by~\citet{mukhoti2020calibrating} (\S3, Fig.~2) and the explicit observation by~\citet{tao2023pcs} that ``early stopping fails to calibrate networks''. What varies across seeds is the epoch at which dev macro-F\textsubscript{1} peaks: four well-calibrated seeds peak at epochs 4--7, two poorly-calibrated seeds peak at epoch 3. Macro-F\textsubscript{1}-based checkpoint selection combined with patience-based early stopping saves the epoch-3 checkpoint for those seeds when dev ECE is still $\sim 0.13$.

Principled remedies in the calibration literature include AdaFocal's validation-driven $\gamma$ adaptation~\citep{ghosh2022adafocal}, predecessor combination (PCS, \citet{tao2023pcs}), and the refine-then-calibrate decomposition~\citep{holzmuller2025refine}. We note an immediately actionable consequence: replacing the pure F\textsubscript{1} selection criterion with a composite score $e^\star = \arg\max_e [\mathrm{F1}(e) - \alpha\,\mathrm{ECE}(e)]$ should collapse the bimodality, since for every ``bad'' seed an alternative epoch exists at which dev ECE is below $0.07$ while dev F\textsubscript{1} is within $0.005$ of its observed peak. The bimodal-ECE observation also explains why post-hoc TS, while improving global ECE, cannot rescue the worst-group calibration on the bad seeds: TS shifts the global scale but not the per-group calibration that breaks under early-stopping truncation.

\subsection{Limitations}

\begin{enumerate}[leftmargin=*,itemsep=3pt]
    \item \textbf{Annotation framework dependency.} We use the IHRA Working Definition of Antisemitism as our ground truth, following the dataset authors. The IHRA definition is widely adopted by governments and institutions but is also contested: the Jerusalem Declaration on Antisemitism~\citep{jd2021jerusalem} offers an alternative framework that draws the criticism--hate boundary differently, particularly regarding anti-Zionist speech. Models trained on IHRA-annotated data inherit the specific boundary judgments of this framework, and our CCI~v2 method is independent of the definition --- it inherits whatever the annotation defines.

    \item \textbf{Single primary dataset, single identity.} Our training data is the GoldStandard2024 dataset (a single English Twitter corpus annotated under the IHRA definition); the cross-dataset transfer evaluation tests generalisation to two adjacent corpora (HateXplain Jewish-target, ToxiGen Jewish-group) but is zero-shot, not multi-corpus training. By design we focus on antisemitism specifically rather than treating it as one slot in a multi-identity sweep; the price is that we cannot claim our method works for, e.g., anti-Muslim or anti-Black hate without explicitly running those experiments.

    \item \textbf{Theoretical scope of the structural claim.} Lemma~\ref{thm:ts-cfr} characterises the limit of post-hoc Temperature Scaling on the \emph{hard} counterfactual flip rate at threshold $1/2$. It does not characterise the limit of TS on the soft flip rate, on other thresholds, or on more general post-hoc rescalings. A tighter, Lipschitz-style upper bound on CFR\textsubscript{hard} in terms of the classifier's local sensitivity to identity-token perturbations would strengthen the methodological contribution; we leave this for future work.

    \item \textbf{Symmetry filter heuristics.} Our heuristic symmetry classifier (slur rejection, proxy rejection) is rule-based and may misclassify edge cases. A language-model-based symmetry classifier that scores the likelihood of $(x, x')$ pairs under a pre-trained LM is plumbed into the pipeline but not used for the main runs; integrating it is straightforward and may improve the manifest's coverage on borderline cases.

    \item \textbf{English-only scope.} Our dataset and models are limited to English-language tweets. Antisemitism manifests differently across languages and cultures; multilingual extension --- in particular cross-language transfer to Hebrew and Arabic --- is a natural follow-up.

    \item \textbf{Binary classification.} We frame antisemitism detection as a binary task, collapsing the rich spectrum of antisemitic expression (overt hate, conspiracy theories, Holocaust trivialization, microaggressions, coded references) into a single positive class. A multi-label or severity-graded approach could provide more actionable information for content moderators.

    \item \textbf{Temporal limitation.} The dataset spans 2019--2023, a period that includes significant geopolitical events. The model's performance on discourse after this period --- with potentially new coded language, geopolitical contexts, and platforms --- is unknown. Continual-learning approaches may be necessary for deployed systems.
\end{enumerate}
